\Askhsh{18781}{4}
\begin{ylh}{1}
    \item 3.1 Ο Κύκλος.
\end{ylh}
Δίνονται δύο κύκλοι με εξισώσεις $C_1: x^2 + y^2 - 6x + 5 = 0$ και $C_2: x^2 + y^2 = 1$.
\begin{alist}
\item Να δείξετε ότι:
    \begin{rlist}
    \item Η εξίσωση του κύκλου $C_1$ γράφεται στη μορφή $(x - 3)^2 + y^2 = 4$. \monades{5}
    \item Οι κύκλοι $C_1, C_2$ εφάπτονται εξωτερικά. \monades{4}
    \end{rlist}
\item Να βρείτε:
    \begin{rlist}
    \item Το σημείο επαφής των δύο κύκλων $C_1$ και $C_2$. \monades{6}
    \item Την εξίσωση της εσωτερικής κοινής εφαπτομένης των δύο κύκλων $C_1$ και $C_2$. \monades{4}
    \end{rlist}
\item Αν τα σημεία $M_1, M_2$ διατρέχουν τους κύκλους $C_1, C_2$ αντίστοιχα, να βρείτε τη μέγιστη απόσταση ανάμεσα στα σημεία αυτά. \monades{6}
\end{alist}